\scanpage{31}
[où il s'agit des polynômes caractéristiques] % L1

Pour ceci, notons que % L2
\[
u\in\mathrm{Hom}(\mathbb 1_X,\mathrm{Hom}(M,\unsure{M?|M_2|M}))=\mathrm{Hom}(\mathbb 1_X,\overset{\vee}{M}\otimes M) % L3
\]
et on a un morphisme \uline{contraction} % L4
\[
\overset{\vee}{M}\otimes M\longrightarrow\mathbb 1_X % L5
\]
% slanted
\margin{\uline{à inclure dans le formalisme tensoriel}} % L5a
d'où un $c(u)\in\mathrm{Hom}(\mathbb 1_X,\mathbb 1_X)$ % L6
et \deleted{\unsure{donc}} $\mathrm{Tr}\,T_\ell(u)=T_\ell(c(u))\in\deleted{\mathbb Z}\mathbb Q_\ell$. % L7
et il suffit de savoir : % L8
\[
\boxed{X\ \text{connexe}\Longrightarrow\mathrm{Hom}(\mathbb 1_X,\unsure{\mathbb 1|1}_X)\simeq\mathbb Q\,\mathrm{id}_{\mathbb 1_X}} % L9
\]

9) \uline{Filtration de} $\mathcal M^+(X)$ \uline{et} $\mathcal M(X)$ % L10

1°) \uline{Filtration par le poids}, \deleted{\unsure{et}} % L11
\[
\deleted{\unsure{\mathcal M^+(X)}}\ \mathcal M^{+0}(X)\subset\mathcal M^{+1}(X)\subset\cdots\subset\mathcal M^{+i}(X)\subset\cdots % L12
\]
\uline{filtration exhaustive} de $\mathcal M^+(X)$. [L'appartenance % L13
à $\mathcal M^{+i}(X)$ se vérifie fibre par fibre, \deleted{\unsure{l'app.}} % L14
% slanted
\margin{indices en bas} % L14a
\inserted{\deleted{\unsure{ou}} (géométrique si on veut)} et dans le cas $X$ de t.f. \inserted{sur $\mathbb Z$} il suffit de % L15
vérifier en les pts fermés. (Pour \unsure{voir|avoir|ceux-} % L16
\unsure{ça|ce|ci}, il y a un critère par Frobenius, cf % L17

% archive sheet number: 30
